\chapter{How to read this document}

\section{What this is}

This is the single reference for the two-point parked-frequency lock-in on the
NV compact magnetometer: how each acquisition method works, exactly where the
acquisition time goes, what sensitivity each method delivers, which defects are
open, and what has been changed in the code so far.

It merges and supersedes two earlier documents:

\begin{itemize}
  \item \file{docs/2026-08-06\_twopoint\_timing/TIMING\_AND\_NOISE\_ANALYSIS.md}
        --- the 2026-08-06 session: timing budget, readout-window optimum,
        discovery of the burst-mode replay.
  \item \file{docs/2026-08-14\_twopoint\_methods/TWOPOINT\_METHODS\_AND\_TIMING.md}
        --- the 2026-08-14 session: the three modes compared, the per-call floor,
        two open defects.
\end{itemize}

Both remain in the repository as the record of what was believed when. Where they
disagree with this document, this document wins, and
Chapter~\ref{ch:retractions} lists every such disagreement explicitly.

\section{Where the numbers come from}

Every measured value here is regenerated by one script from committed CSV files:

\begin{center}
\co{python scripts/analyze\_twopoint\_0814.py}
\end{center}

It writes \file{docs/2026-08-14\_twopoint\_methods/tables/*.csv} and
\file{figures/*.png}; this document quotes those tables. Nothing is hand-typed,
so re-running the script after a new session refreshes the numbers rather than
invalidating them. Values that can only be measured on the rig are marked
\Open{} and named as such --- there are five of them, all in
Chapter~\ref{ch:todo}.

\section{The three findings that matter most}

\begin{keybox}[1. The acquisition rate is set by the host, not by the readout window]
An \co{acquire()} call costs \textbf{10--11.4\,ms of host round-trip whatever the
program does}, and the FPGA pulse work runs \emph{underneath} that rather than
adding to it. Halving the readout window from \SI{213}{\micro\second} to
\SI{120}{\micro\second} cut the FPGA work from 9.80\,ms to 2.40\,ms and changed
the rate by nothing (87.5 $\rightarrow$ 86.8\,Hz).

Consequences: shortening $\tau$ is not a rate lever in averaged mode; averaging
more reps is \emph{free} up to the call floor; and beating $\sim$90\,Hz requires
amortising the call across many samples, which is what burst and streaming do.
See Chapter~\ref{ch:timing}.
\end{keybox}

\begin{keybox}[2. The photodiode fix bought 4.1$\times$, and burst mode shows it]
After the reference-photodiode gain was corrected and the signal diode stopped
saturating, the burst-mode noise floor went from \textbf{158 to 38\,\etaunit} ---
a factor of 4.1 at the same two parked points. Burst is the best operating point
available today by a wide margin. See Chapter~\ref{ch:sensitivity}.
\end{keybox}

\begin{keybox}[3. Streaming reaches 1\,kHz, but its noise floor is $\sim$6$\times$ burst]
Streaming delivered \textbf{1041.7\,Hz} with a gapless, duplicate-free rep axis
across 125\,000 reps --- the 1\,kHz target is met. Its noise floor, however, sits
a flat $\sim$6$\times$ above burst across 10--500\,Hz for reasons not yet
identified. Chapter~\ref{ch:defects} lists what has been ruled out and the
one-minute rig test that decides it.
\end{keybox}

\section{Status at a glance}

\begin{center}
\small
\begin{tabular}{@{}llll@{}}
\toprule
\hd{Item} & \hd{State} & \hd{Number} & \hd{Where} \\
\midrule
Averaged mode rate            & \Done    & 86.8--88.8\,Hz          & \S\ref{sec:avg} \\
Burst mode rate               & \Done    & 4167\,Hz in-burst       & \S\ref{sec:burst} \\
Streaming rate                & \Done    & 1041.7\,Hz              & \S\ref{sec:stream} \\
Best measured noise floor     & \Done    & 38\,\etaunit\ (burst)   & \S\ref{sec:etameas} \\
Readout window chosen         & \Done    & \SI{120}{\micro\second} & \S\ref{sec:window} \\
Timing model                  & \Done    & $\max(\text{host},\text{FPGA})$ & \S\ref{sec:model} \\
Free averaging implemented    & \Done    & 42 reps fit, 10 used    & \S\ref{sec:freeavg} \\
Per-mode post-processing      & \Done    & 3 pipelines             & \S\ref{sec:postproc} \\
Notebook reorganised          & \Done    & 4A/4B/4C + 5A/5B/5C     & \S\ref{sec:notebook} \\
\midrule
Burst replay root cause       & \Open    & 49.9\% of batches       & \S\ref{sec:replay} \\
Streaming noise excess        & \Open    & $\sim$6$\times$ burst   & \S\ref{sec:streamnoise} \\
Per-acquire noise term        & \Open    & does not average down   & \S\ref{sec:notshot} \\
ABBA ordering benefit         & \Open    & never tested on hardware & \S\ref{sec:abba} \\
Calibrated system sensitivity & \Open    & no field reference yet  & \S\ref{sec:notaclaim} \\
\bottomrule
\end{tabular}
\end{center}


\chapter{The measurement, from photons to \texorpdfstring{$\Delta f$}{delta f}}
\label{ch:estimator}

Before any timing discussion, this is what a ``sample'' actually is. All three
acquisition methods produce the same quantity by the same arithmetic; they differ
only in how many FPGA repetitions go into one sample and how the host collects
them.

\section{Two parked frequencies on one flank pair}

A single ODMR resonance is fitted (Lorentzian, from the Step~1 reference sweep),
giving its centre $f_0$, half-width $\gamma$, baseline $B$, and the two
maximum-slope points $f_-$ and $f_+$ on either flank. Rather than sweeping, the
microwave source is \emph{parked} alternately at $f_-$ and $f_+$ and the
photoluminescence is read at each.

\begin{align}
z_\pm &= \frac{\text{counts at } f_\pm}{\text{normalisation}}
   &&\text{normalised PL at each parked point} \\
L &= z_+ - z_-  &&\text{the lock-in signal} \\
\Delta f &= \frac{L - (b_+ - b_-)}{m_- - m_+}
   &&\text{shift of the resonance, in MHz}
\end{align}

where $b_\pm$ are the calibration baselines at the parked points and $m_\pm$ the
flank slopes, both measured off the reference sweep. Field is then
$\Delta B = \Delta f / \gamma_\text{NV}$ with
$\gamma_\text{NV} = \SI{28.024}{\kilo\hertz\per\micro\tesla}$ along the tracked
NV axis.

\subsection{Why the difference and not one flank}

If the resonance shifts by $\delta$, $z_-$ moves down the negative-slope flank
and $z_+$ moves up the positive-slope one: the shift appears with opposite sign
in the two channels and \emph{adds} in $L$. A change in laser power or collection
efficiency scales both channels the same way and \emph{cancels} in $L$ to first
order.

This works, and by a large factor. In the 2026-08-06 heat-gun/motor run, $z_-$
and $z_+$ swung \textbf{13.4\% peak-to-peak together} while the difference stayed
flat --- a common-mode rejection ratio of \textbf{32$\times$}.

\subsection{The noise consequence}

Because $\Delta f$ is linear in $L$,
\begin{equation}
\sigma(\Delta f) = \frac{\sigma(z_+ - z_-)}{|m_- - m_+|}.
\label{eq:noise}
\end{equation}
Two things follow that recur throughout this document:

\begin{enumerate}
  \item \textbf{Collapsed dip depth costs sensitivity directly.} The slopes come
  from the calibration. If the resonance moves off the parked pair, or the
  laser/MW level drifts, the real slopes flatten while the assumed ones do not,
  and $\sigma(\Delta f)$ grows in proportion. Every acquisition cell now prints
  the measured dip depth at both parked points against the calibrated value.
  \item \textbf{Post-processing must act on the difference, not on the channels.}
  Filtering $z_-$ and $z_+$ independently breaks the cancellation. This is not
  hypothetical: see \S\ref{sec:despikebug}.
\end{enumerate}

\section{The residual first-order term, and ABBA}
\label{sec:abba}

The cancellation is not perfect, because the two points are not read at the same
instant. In the default \co{"forward"} ordering the rep visits $f_-$ then $f_+$,
one readout window apart, \emph{always in the same direction}. A common-mode
drift $d$ per unit time therefore contributes $+d\,\tau$ to $L$ every rep: the
estimator is differentiating the common mode.

\co{"abba"} ordering emits $f_-, f_+, f_+, f_-$ and folds slot 0 with slot 3 and
slot 1 with slot 2. A linear drift contributes $+d, +2d, +3d, +4d$ across the
four slots, and the folding cancels it exactly. It doubles the readouts per rep,
so one \co{"abba"} rep costs and averages exactly two \co{"forward"} reps --- the
cancellation is paid for in rep count, not in time.

\begin{openbox}[OPEN --- ABBA has never been tested on hardware]
Implemented and verified in simulation; free at equal averaging. Given the
32$\times$ rejection already measured and the 13.4\% common-mode swing, the
residual first-order term is plausibly significant, but no A/B has been run.
Pass criterion and protocol: \S\ref{sec:rigchecks}, check S7.
\end{openbox}

\section{The two sensitivity conventions, fixed once}
\label{sec:conventions}

Two numbers are quoted throughout, and they are not the same thing.

\begin{description}[style=nextline,leftmargin=1.2em]
  \item[$\eta = \sigma_B\sqrt{2\,\Delta t}$ --- the ``implied'' sensitivity]
  The white-noise amplitude spectral density that a measured standard deviation
  $\sigma_B$ at sample spacing $\Delta t$ would correspond to. Cheap to compute
  from any trace. \emph{Only meaningful if the trace is actually white}: drift
  inflates $\sigma$ without raising the noise density.
  \item[Measured floor --- the median ASD in the flat band]
  Read directly off a Welch periodogram between 10\,Hz and $0.4 f_s$. Makes no
  whiteness assumption.
\end{description}

Both are reported side by side, and \textbf{when they disagree the floor is the
one to trust}; the ratio is a direct readout of how drift-dominated the trace is.
For example the 2026-08-14 burst run gives $\eta = 56$ but a floor of
38\,\etaunit, a ratio of 1.5 --- the trace is not white.

The convention is fixed in code, in \file{notebook_modules/twopoint_spectra.py},
whose PSD normalisation is verified against Parseval's theorem (the integral of
the PSD reproduces the variance to 0.02\%) and against a synthetic tone of known
amplitude.

\begin{retractbox}[Correction --- an earlier $\sqrt{2}$ and an earlier $\sqrt{\text{reps}}$]
The 2026-08-06 document used $\eta = \sigma_B\sqrt{\Delta t}$, smaller than the
convention above by $\sqrt{2}$. Its figures are restated here in the
$\sqrt{2\Delta t}$ convention where they are compared with 2026-08-14 numbers.

Separately, the first draft of that document paired a \emph{reps-averaged}
sweep-point $\sigma$ with a \emph{single-rep} time, understating $\eta$ by
$\sqrt{\text{reps}}\approx 10$. That was corrected on 2026-08-12 by
cross-calibrating against the burst runs; see \S\ref{sec:window}.
\end{retractbox}

\section{What is not being claimed}
\label{sec:notaclaim}

Every sensitivity figure in this document is a \textbf{readout-chain} figure: it
comes from the measured scatter of $\Delta f$ and the ODMR-derived slopes. None
of them is a calibrated system sensitivity, because that requires a known applied
field, which has not been done. The distinction matters for the project's
sensitivity claims and is the reason
\file{NV Project/wiki/concepts/sensitivity.md} keeps these numbers separate from
the roadmap target.
